% =====================================================================
\chapter{Словарь терминов}
% =====================================================================

\begin{small}
\renewcommand{\arraystretch}{1.15}
\begin{longtable}{@{}L{40mm}L{120mm}@{}}
\toprule
\textbf{Термин} & \textbf{Значение} \\
\midrule
\endhead
Aileron / Roll & элероны / крен: наклон модели на бок \\
Elevator / Pitch & руль высоты / тангаж: нос вверх--вниз \\
Rudder / Yaw & руль направления / рыскание: поворот носа влево--вправо \\
Throttle & газ: скорость моторов \\
AETR / TAER & порядок первых четырёх каналов: Aileron, Elevator, Throttle, Rudder \\
AUX & вспомогательные каналы (\sw{CH5} = AUX1, \sw{CH6} = AUX2 …) \\
Arm (арм) & разрешение работы моторов в полётном контроллере \\
Binding (привязка) & связывание передатчика с конкретным приёмником \\
Binding phrase & фраза привязки в ExpressLRS \\
Bootloader & загрузчик; программа, которая записывает прошивку \\
CRSF & протокол Crossfire, используется между пультом и модулем ELRS и между приёмником и
  контроллером \\
Crow (кроу) & тормозная конфигурация: элероны вверх, закрылки вниз \\
DFU & режим прошивки микроконтроллера STM32 по USB \\
Dual rate (расходы) & переключаемый максимальный ход управления \\
ELRS (ExpressLRS) & открытая радиосистема управления \\
Elevon (элевон) & поверхность летающего крыла, совмещающая элерон и руль высоты \\
Expo & экспонента; смягчение управления вблизи центра стика \\
Failsafe & поведение модели при потере связи \\
Flight mode (FM) & полётный режим EdgeTX: набор триммеров и условий \\
GVar (GV) & глобальная переменная модели \\
Hall gimbal & стик на датчиках Холла (бесконтактный, без износа) \\
Input & вход; обработанный сигнал стика («пилотская» часть) \\
JR bay & отсек для внешнего модуля стандартного размера \\
LQ (Link Quality), RQly & качество связи: доля принятых пакетов, \% \\
Mode 1/2 & раскладка стиков; Mode 2 --- газ слева \\
Logical switch & логический (виртуальный) переключатель \\
Lua & встраиваемый язык скриптов EdgeTX \\
Mix & миксер; правило, как входы попадают в каналы \\
Model Match & защита от управления «не той» моделью \\
Multi (4-в-1) & мультипротокольный модуль \\
Output & выход; «механическая» настройка канала \\
Packet rate & частота отправки пакетов управления, Гц \\
PPM & аналоговый формат передачи нескольких каналов по одному проводу \\
RSSI & уровень принимаемого сигнала, дБм \\
Special function & специальная функция: действие по условию \\
Subtrim & сдвиг нейтрали канала на выходе \\
Telemetry & данные, передаваемые от модели на пульт \\
Trainer & режим «учитель--ученик» \\
Trim & триммер; оперативная подстройка нейтрали в полёте \\
Сервомашинка (серва) & моторчик с рычагом, который двигает руль модели \\
Полётный контроллер (FC) & компьютер на борту коптера, который крутит моторы \\
YAML & текстовый формат файлов моделей и настроек EdgeTX \\
\bottomrule
\end{longtable}
\end{small}
